% !PRINTED 505
\weberpage{505}

%<<CONT
\gesperrt{durch keine andere Operation aus $P$ in eine bestimmte von $\psi$ verschiedene Grösse $\psi_1$ über.}

Demnach entsprechen den $j$ Nebengruppen ebenso viele Functionen $\psi$, nämlich
\weq{10}{\psi,\ \psi_1,\ \psi_2\ldots\psi_{j-1},}
und aus dem Theorem 2. ergiebt sich der Satz:

5. \gesperrt{Die Grössen (10) erleiden eine Permutation, wenn auf alle gleichzeitig eine und dieselbe Operation $\pi$ aus $P$ angewandt wird.}

Denn wir können die Grössen (10) so darstellen
\[
\psi(\varkappa),\ \psi(\varkappa\pi_1),\ \psi(\varkappa\pi_2),\ldots \psi(\varkappa\pi_{j-1}),
\]
und wenn also $\pi$ darauf angewandt wird,
\[
\psi(\varkappa\pi),\ \psi(\varkappa\pi_1\pi),\ \psi(\varkappa\pi_2\pi),\ldots \psi(\varkappa\pi_{j-1}\pi),
\]
worin $\varkappa$ und $\pi$ beliebige Permutationen aus $Q$ und $P$ sind. In jeder der beiden Reihen
\[
\begin{aligned}
&\varkappa,\ \varkappa\pi_1,\ \varkappa\pi_2\ldots\varkappa\pi_{j-1}\\
&\varkappa\pi,\ \varkappa\pi_1\pi,\ \varkappa\pi_2\pi\ldots\varkappa\pi_{j-1}\pi
\end{aligned}
\]
kommt aber aus jeder der Nebengruppen (4) eine und nur eine Permutation vor.

Die Functionen (10), die, wie wir gesehen haben, alle von einander verschieden sind, heissen \gesperrt{conjugirte Functionen}.

Wir können nach dem Vorhergehenden leicht die Gruppen bilden, zu denen die einzelnen Functionen (10) gehören. Es ist z.\,B. $\psi_1 = \psi(\pi_1)$, und wenn nun $\pi$ eine Permutation aus $P$ ist, durch die $\psi_1$ ungeändert bleibt, so muss
\[
\psi(\pi_1\pi) = \psi(\pi_1)
\]
sein. Auf diese Gleichung können wir aber die in $P$ enthaltene Permutation $\pi_1^{-1}$ anwenden, wodurch sich
\[
\psi(\pi_1\pi\pi_1^{-1}) = \psi
\]
ergiebt; d.\,h. es muss $\pi_1\pi\pi_1^{-1}$ zu $Q$ gehören. Setzen wir es gleich $\varkappa$, so folgt
\weq{11}{\pi = \pi_1^{-1}\varkappa\pi_1.}
Wenn umgekehrt $\pi$ diese Form hat, so ist
\[
\psi_1(\pi) = \psi(\pi_1\pi) = \psi(\pi_1\pi_1^{-1}\varkappa\pi_1) = \psi(\varkappa\pi_1) = \psi_1,
\]
und $\psi_1$ gestattet also diese Permutation. Demnach können wir die Gruppe, zu der $\psi_1$ gehört, durch das Symbol
\[
\pi_1^{-1}Q\pi_1
\]
%CONT>>
